\begin{enumerate}[leftmargin=2.5em]
\item[] \printedpage{34}%
% head: A. Grothendieck : Dérivateurs, Ch. XVI, Édition des Universités de Montpellier II et Paris VII
\textit{the rectangle of implications}
\[
\begin{tikzcd}[row sep=large, column sep=large, execute at end picture={
\draw[decorate,decoration={brace,amplitude=6pt}] ([xshift=-14pt]A.west |- C.south) -- ([xshift=-14pt]A.north west) node[midway,left=8pt]{$\mathcal{F}_W,\mathrm{Fibf}_W$};
\draw[decorate,decoration={brace,amplitude=6pt}] ([xshift=14pt]AW.north east) -- ([xshift=14pt]AW.east |- CW.south) node[midway,right=8pt]{$\mathcal{F}_{W_0},\mathrm{Fibf}_{W_0}$};
}]
|[alias=A]| \mathrm{A} \arrow[r,Rightarrow] \arrow[d,Rightarrow] & |[alias=AW]| \mathrm{A}_{W_0} \arrow[d,Rightarrow] \\
\mathrm{B} \arrow[r,Rightarrow] \arrow[d,Rightarrow] & \mathrm{B}_{W_0} \arrow[d,Rightarrow] \\
|[alias=C]| \mathrm{C} \arrow[r,Rightarrow] \arrow[d,Rightarrow] & |[alias=CW]| \mathrm{C}_{W_0} \arrow[d,Rightarrow] \\
\underbrace{\mathrm{D}}_{\mathrm{Parf}_W} \arrow[r,Rightarrow] & \underbrace{\mathrm{D}_{W_0}}_{\mathrm{Parf}_{W_0}}\,,
\end{tikzcd}
\tag*{(4.4.13.5)}
\]
\textit{where we have recalled, next to the braces, what the set, or the sets $\mathcal{F}\subset\mathrm{Fl}(\mathrm{Cat})$ are which define the conditions considered (as one of the conditions $\mathrm{A}(\mathcal{F})$, $\mathrm{B}(\mathcal{F})$, $\mathrm{C}(\mathcal{F})$ or $\mathrm{D}(\mathcal{F})=\mathrm{A}(\mathcal{F})$ in the case where the three conditions are equivalent). The weakest of the eight conditions considered is $\mathrm{D}_{W_0}$, the strongest} A. (\textit{And for $W_0$ variable, we find the weakest condition of all by taking $D_{W_\omega}$, where $W_\omega$ is the smallest of the basic localizers in} Cat.)

\item[c)] \textit{We have moreover, for the weakest condition $\mathrm{D}_{W_0}$, the implication}
\[
\mathrm{D}_{W_0} \implies \mathrm{F},
\]
\msp{page 60}%
\textit{where} $\mathrm{F}=\text{Loc (7 F)}$ \textit{is the condition that} Cat \textit{has enough $W$-fibrations,} i.e. \textit{that it is a category of fibrations for} $(W,\mathrm{Fib}_W)$.
\end{enumerate}

\noindent\cercle{2}\enspace{\itshape If $W$ and $W_0$ satisfy \textup{Loc (8 bis)} (axiom of countable limits), then \textup{$\mathrm{D}_{W_0}\Longrightarrow\mathrm{A}$, i.e.} the 8 conditions considered are equivalent. Thus if $W$ satisfies \textup{Loc (8 bis)}, considering the smallest basic localizer $W_L$ satisfying \textup{Loc (8 bis)} (so that $W\supset W_L$), we see that the eight conditions
\[
\mathrm{A},\ \mathrm{B},\ \mathrm{C},\ \mathrm{D},\ \mathrm{A}_{W_L},\ \mathrm{B}_{W_L},\ \mathrm{C}_{W_L},\ \mathrm{D}_{W_L}
\]
are equivalent~\footnote[33]{N.B. The weakest form, $\mathrm{D}_{W_{\mathrm{L}}}$, of these conditions, is the relation \fbox{$\mathrm{Parf}_{W_{\mathrm{L}}}\cap W\subset W^{\mathrm{u}}$.}}. In particular, under the same (quite innocuous) hypothesis on $W$, the conditions \textup{Loc (7) $=$ Loc (7 a), Loc (7 b)} and \textup{Loc (7 c)} are equivalent, and they are also equivalent to the condition \textup{$\mathrm{D}=\mathrm{A}(\mathrm{Parf}_W)$, i.e.} to
\[
\mathrm{Parf}_W\subset\mathrm{Fib}_W \qquad (\textit{special case of } \text{Loc (7 a)}),
\]
and to its two seemingly weaker variants \textup{$\mathrm{B}(\mathrm{Parf}_W)$, $\mathrm{C}(\mathrm{Parf}_W)$} (special case of \textup{Loc (7 b) and of Loc (7 c)} respectively).}

% foot: Version 18 mars 2019

\noindent\cercle{3}\enspace{\itshape Suppose that $W_0$ satisfies \textup{Loc (8 bis)} (without any particular hypothesis on $W$). Then
\printedpage{35}%
% head: A. Grothendieck : Dérivateurs, Ch. XVI, Édition des Universités de Montpellier II et Paris VII
\msp{page 61}%
the four conditions $\mathrm{A}_{W_0}$, $\mathrm{B}_{W_0}$, $\mathrm{C}_{W_0}$, $\mathrm{D}_{W_0}$ are equivalent, and the diagram of implications of \textup{1°)} (p. \textup{59}) therefore reduces to the linear diagram
\[
\mathrm{A} \implies \mathrm{B} \implies \mathrm{C} \implies \mathrm{D} \implies \mathrm{D}_{W_0}\ (\Longleftrightarrow \mathrm{A}_{W_0}\ldots),
\]
whose first term is \textup{Loc (7 a) $=$ Loc (7)}, and the last is equivalent to the inclusion
\[
\mathcal{F}_{W_0}\subset\mathrm{Fib}_W,
\]
special case of \textup{$\mathrm{A}=$ Loc (7 a)}, or again to the inclusion (apparently stronger)
\[
\mathrm{Fibf}_{W_0}\subset\mathrm{Fib}_W,
\]
(special case of the equivalent form
\[
\mathrm{Fibf}_W\subset\mathrm{Fib}_W,
\]
of the axiom} Loc (6 a)).

\noindent\cercle{4}\enspace\textit{Finally, we have the following more technical conditional implications.}
\begin{enumerate}[leftmargin=2.5em]
\item[a)] \textit{If $W_0$ satisfies Loc (7 F), then the conditions $\mathrm{A}_{W_0}$, $\mathrm{B}_{W_0}$, $\mathrm{C}_{W_0}$ are equivalent, so diagram (4.4.13.5) reduces to the diagram}
\[
\begin{tikzcd}[row sep=large, column sep=large]
\mathrm{A} \arrow[d,Rightarrow] & \\
\mathrm{B} \arrow[d,Rightarrow] & \\
\mathrm{C} \arrow[r,Rightarrow] \arrow[d,Rightarrow] & \mathrm{C}_{W_0} \arrow[d,Rightarrow] \\
\mathrm{D} \arrow[r,Rightarrow] & \mathrm{D}_{W_0}\ .
\end{tikzcd}
\]

\item[b)] \textit{If $W$ is strongly saturated} (\textit{axiom} Loc (1 ter), \textit{which}
\msp{page 62}%
\textit{is satisfied in the great majority of cases, and is even imperative for constructing a prederivator satisfying the localization axiom} Der 2) \textit{and satisfies} Loc (3 bis) (\textit{strong localization axiom,} cf. §1), \textit{then we have} $\mathrm{B}\Longrightarrow\mathrm{A}$, \textit{hence}
\[
\mathrm{A} \iff \mathrm{B},
\]
i.e. Loc (7 a) $=$ Loc (7) \textit{is equivalent to} Loc (7 b).
\end{enumerate}

\noindent\textsc{Proof}. Everything stated has been seen in the preceding pages, with the sole exception of 4°, b. But this implication $\mathrm{B}\Longrightarrow\mathrm{A}$ for $W$ strongly saturated is proved in XII (prop. 1, p. 54). I recall that the proof makes essential use of the % [sic] prouvé -> prouvée, accord avec « implication »
% foot: Version 18 mars 2019
\printedpage{36}%
% head: A. Grothendieck : Dérivateurs, Ch. XVI, Édition des Universités de Montpellier II et Paris VII
formalism of the prederivator $\mathrm{HOT}_W$ (constructed and studied mainly in VI), and the theory of the functors $f_!$, and in particular the result (valid if $W$ satisfies Loc (3 bis, 7 b)) that for an arrow $f:S'\longrightarrow S$ in $W$, the corresponding arrow
\[
f^*_{\mathrm{lc}}:\mathrm{HOT}^{\mathrm{lc}}_W(S)\longrightarrow\mathrm{HOT}^{\mathrm{lc}}_W(S')
\]
is an equivalence of categories. So the implication $\mathrm{B}\Longrightarrow\mathrm{A}$, too, is a result which is not at all obvious.

\noindent\msp{page 63}%
4.4.14. I summarize the implications obtained so far, including the conditional implications, in the following summary diagram:
\[
\begin{tikzpicture}[baseline=(C.base),font=\small]
\node[draw,circle,inner sep=1pt] (A) at (0,0) {$\mathrm A$};
\node (AW) at (2,0) {$\mathrm A_{W_0}$};
\node (B) at (0,-1.2) {$\mathrm B$};
\node (BW) at (2,-1.2) {$\mathrm B_{W_0}$};
\node (C) at (0,-2.4) {$\mathrm C$};
\node (CW) at (2,-2.4) {$\mathrm C_{W_0}$};
\node (D) at (0,-3.6) {$\mathrm D$};
\node (DW) at (2,-3.6) {$\mathrm D_{W_0}$};
\node (F) at (7.8,-3.6) {$\mathrm F$};
\node[draw,circle,inner sep=1pt] (Ab) at (2,-4.8) {$\mathrm A$};
\node[anchor=west] at (2.25,-4.8) {$\,,$};
\foreach \p/\q in {A/AW,B/BW,C/CW,D/DW,A/B,B/C,C/D,AW/BW,BW/CW,CW/DW,DW/F}
  \draw[double,double distance=1pt,-{Implies}] (\p) -- (\q);
\draw[dotted,->] (B.west) to[out=145,in=215] node[left,font=\scriptsize] {$W$ (1 ter, 3 bis)} (A.west);
\draw[dotted,->] (CW.east) to[out=70,in=-70] node[right,font=\scriptsize] {$W_0$ (7)} (AW.east);
\draw[dotted,->] (DW.east) to[out=60,in=-40,looseness=1.45] node[right,font=\scriptsize] {$W_0$ (8 bis)} (AW.east);
\draw[dotted,->] (D.south east) -- node[below left,font=\scriptsize] {$W$ (8 bis)} (Ab.north west);
\draw[dotted,->] (DW) -- node[right,font=\scriptsize] {$W$ (8 bis) $+$ $W_0$ (8 bis)} (Ab);
\node[anchor=south west,align=left,font=\scriptsize] at (3.4,-3.45) {N.B. we shall see from the example of $\mathbf{W}_0$\\
that this implication is \textit{strict,}\\
even under the hypothesis\\
Loc (1 ter, 2 bis, 4, 5, 6 a, 8, 9).};
\end{tikzpicture}
\]
where I have circled the strongest and most important condition of all, $\mathrm{A}=$ Loc (7), which appears at two opposite ends of the diagram.

I would like to make one more comment on the ingredients of the results proved here, and in particular of the equivalence of the eight conditions A to $\mathrm{D}_{W_0}$, when $W$ and $W_0$ satisfy Loc (8 bis):

\msp{page 64}%
\begin{enumerate}[leftmargin=2.5em]
\item[1)] The theory of \textit{categories of paths} \underline{Ch}$(X)$ (VII), relying on that of \textit{smooth, proper, perfect morphisms.}
\item[2)] The lemma of K. S. \textsc{Brown} (3.2.1, page 20).
\item[3)] For the implication $\mathrm{B}\implies\mathrm{A}$, \textit{i.e.} Loc (7 b) $\implies$ Loc (7) (under the hypothesis $W$ (1 ter, 3 bis)), I use the theory of the prederivator $\mathrm{Hot}_W$ (VI), and more particularly that of the functors $f_!$ (defined under the sole hypothesis $W$ (3 bis), strengthening the automatic condition $W$ (3) (forming part of the definition of a basic localizer)).
\end{enumerate}
% foot: Version 18 mars 2019

\leavevmode\printedpage{37}%
% head: A. Grothendieck : Dérivateurs, Ch. XVI, Édition des Universités de Montpellier II et Paris VII
It is clear that I would not have managed to see clearly, without the fundamental work of K. S. \textsc{Brown}. Yet there still remain open questions. Thus, if we do not assume that $W$ satisfies Loc (8 bis), but if need be Loc (3 ter) (and I have not yet had occasion to work with a localizer which does not satisfy it \ldots), is it nevertheless true that we still have the implication
\[
\mathrm{D} \qquad \Longrightarrow \qquad \mathrm{A},
\]
so that the conditions A, B, C, D are equivalent (and in particular Loc (7 c) already implies Loc (7) in all its force)? It is clear in all this that
\msp{page 65}%
I have not applied myself enough to developing examples and counterexamples. Not to mention that it seems to me that all the basic localizers on Cat which we encounter ``in practice'' already satisfy Loc (8) and a fortiori Loc (8 bis), so for them we do indeed have the expected implication. The same remark applies to the hypothetical implication
\[
\mathrm{D}_{W_0} \overset{?}{\implies} \mathrm{A},
\]
under hypothesis Loc (8 bis) on $W_0$ and on $W$. An obvious candidate for a counterexample already to the first question would be $W_\omega=$ smallest basic localizer on Cat, one that it is almost obvious that it does \textit{not} satisfy Loc (8) nor even Loc (8 bis). But there is no reason why it should satisfy F either, so that (if it does not satisfy it, hence not D either) the implication $\mathrm{D}\Longrightarrow\mathrm{A}$ is indeed valid. This brings us to the question, a more substantial one, whether we might not have the implication % [sic] un qu'il est -> dont il est (?), construction défectueuse
\msp{page 66}%
\[
\mathrm{F} \overset{?}{\implies} \mathrm{A}\ \text{\footnotemark[34]},
\]
\footnotetext[34]{This is false, \textit{cf.} below, the implication Loc (7 $\mathrm{D}_{W_\omega}$) $\Longrightarrow$ Loc (7 F) is \textit{strict.}}%
at least under suitable fairly innocuous conditions on $W$, of the sort Loc (1 ter, 3 bis, 8 bis).

This is probably the most substantial and the most interesting question which is raised by the diagram on page 63, and which no longer calls for the introduction of an auxiliary basic localizer $W_0$. In view of Brown's lemma, it does not seem to me absolutely ruled out that such a miracle should indeed take place!

4.4.15 (Digression on the axiom Loc (6) of cofibrations.) This draws my attention to the in some sense dual question: If Cat has enough $W$-\textit{co}fibrations, \textit{i.e.} if equipped with $(W,\mathrm{Cof}_W)$ it is a category with cofibrations, does this already imply the axiom Loc (6 a) (which had been isolated, essentially, to ensure that this is indeed so). Proposition 3.6. (p. 25) tells us that it amounts to the same thing to prove that, if $W$ satisfies Cof, then the open M. V. squares are
\msp{page 67}%
$W$-cocartesian. But it seems to me now that we can prove this, under the sole additional condition Loc (3 bis)~\footnote[35]{and perhaps, in addition, Loc (1 ter), \textit{i.e.} $W$ strongly saturated.}, which ensures the existence of the functors $f_!$, and % [sic] cocartésien -> cocartésiens, accord avec « les carrés » ; [sic] (note 35) saturée -> saturé, accord avec « W » (localiseur)
% foot: Version 18 mars 2019
\printedpage{38}%
% head: A. Grothendieck : Dérivateurs, Ch. XVI, Édition des Universités de Montpellier II et Paris VII
in particular the construction of Holim of a diagram $\fourche$ \quad as $\int\fourche$ \quad .
So to say that an open $\mathcal{M}$ diagram $\carre$ is $W$-cocartesian, must amount to saying that the canonical arrow
\[
\int\fourche \qquad \longrightarrow \qquad X
\]
is in $W$; at least if $W$ is assumed strongly saturated. Now this has indeed been seen in XII 3.4 (theorem 4 and its \textit{corollary} 3, p. 149). In loc. cit. I state the conclusions under the hypothesis which I called W (7 bis) (the present Loc (6)), but I have just checked that for its part \textit{in the forward direction,} hence also for its corollaries 1 to 3, we do not need
\msp{page 68}%
this hypothesis (as I indicate in marginal annotations).

4.16. The result I have just proved (without carefully checking the proof) seems to me as delicate, as it is satisfying. This makes the hesitant ``conjecture'' of 4.4.14, namely that
\[
\mathrm{F} \overset{?}{\Longrightarrow} \mathrm{A} \qquad \text{hence} \qquad \mathrm{F}\Longleftrightarrow\mathrm{A}
\]
assuming that $W$ satisfies Loc (1 ter, 3 bis, 8 bis), decidedly more inviting. For one begins to get used to the most astounding miracles \ldots

Yet I have just checked that this conjecture is false, using the example of
\[
\mathbf{W}_0=\{f\in\mathrm{Fl}(\mathrm{Cat})\mid\pi_0(f)\ \mathrm{iso}\}\,,
\]
which satisfies F, but not A, nor even $\mathrm{D}_{W_\omega}$. Thus the implication
\[
\mathrm{D}_{W_\omega} \implies \mathrm{F}
\]
and a fortiori the implications
\[
\left\{
\begin{aligned}
&\text{Loc (7 a)} \implies \mathrm{F},\ \text{Loc (7 b)} \implies \mathrm{F},\ \text{Loc (7 c)} \implies \mathrm{F}\\
&\mathrm{D}_{W_0} \implies \mathrm{F}
\end{aligned}
\right.
\]
are \textit{strict,} even when restricting ourselves to basic localizers $W$ satisfying (like $\mathbf{W}_0$)
\[
\text{Loc (1 ter, 2, 3 bis, 4, 5, 6 a, 8, 9)}\ \text{\footnotemark[36]}.
\]
\footnotetext[36]{\textbf{N.B. $\mathbf{W}_0$} does \textit{not} satisfy Loc (6 b), nor Loc (7 $\mathrm{D}_{W_\omega}$) (\textit{cf.} below), on the other hand Loc (7 F).}
% foot: Version 18 mars 2019